%%
%% Track 2 — AI Disclosure Statement Template
%% CHIWork 2026 Workshop:
%%   "Interrogating GenAI Augmentation for CHIworkers:
%%    Strategies for Professional Autonomy and Accountability"
%%
%% Workshop website: https://chiwork-ai.hauke.haus/
%% Workshop day: June 22, 2026 — Linz, Austria
%%
%% INSTRUCTIONS
%% ─────────────────────────────────────────────────────────────
%% A short paragraph is all that is required.
%% This template provides a structured framework if you want
%% to go deeper — use as much or as little as you like.
%%
%% This template is optional. You may submit a plain PDF,
%% a Word document, or any readable format.
%% ─────────────────────────────────────────────────────────────

\documentclass[sigconf]{acmart}

\settopmatter{printacmref=false}
\renewcommand\footnotetextcopyrightpermission[1]{}
\pagestyle{plain}

\acmConference[CHIWork '26]{CHIWork '26 Workshop: Interrogating GenAI Augmentation for CHIworkers}{June 22, 2026}{Linz, Austria}
\acmYear{2026}

%% ── Formatting helpers ────────────────────────────────────────
\usepackage{xcolor}
\definecolor{workshopaccent}{HTML}{6B2737}
\usepackage{tikz}
\usepackage{fontawesome5}
\newcommand{\sectionrule}{\vspace{0.3em}\noindent\textcolor{workshopaccent}{\rule{\linewidth}{0.4pt}}\vspace{0.3em}}
\newcommand{\prompt}[1]{\vspace{0.5em}\noindent\textcolor{workshopaccent}{\small\sffamily #1}\vspace{0.3em}}
\newcommand{\response}{\vspace{0.3em}\noindent\textit{[Your response here.]}\vspace{0.5em}}
\newcommand{\phaseicon}[1]{\textcolor{workshopaccent}{#1}\enspace}

%% ── Your metadata ─────────────────────────────────────────────
\title{AI Use Disclosure: [Your Role / Context / Field]}
%% Examples:
%%   "AI Use Disclosure: A UX Researcher at a Mid-Size Tech Company"
%%   "AI Use Disclosure: PhD Student in HCI, Qualitative Methods Focus"

\author{Your Name}
\affiliation{%
  \institution{Your Institution or Company}
  \city{City}
  \country{Country}}
\email{your@email.com}
%% Pseudonymous submissions are welcome.

%% ─────────────────────────────────────────────────────────────
\begin{document}
\maketitle

%% ─────────────────────────────────────────────────────────────
%% INTRO NOTE — Printed in the PDF for the reader's context
%% ─────────────────────────────────────────────────────────────
\begin{quote}
\small\itshape
This disclosure was submitted to the CHIWork 2026 workshop
\emph{Interrogating GenAI Augmentation for CHIworkers}~\cite{sandhaus2026interrogating}.
We ask: What do you use AI for? What feels right? What feels wrong?
When do you stop? A short paragraph is enough; the framework below
is there if you want to go deeper.
\end{quote}

\sectionrule

%% ═════════════════════════════════════════════════════════════
%% OPTION A — SHORT PARAGRAPH (minimum requirement)
%% ═════════════════════════════════════════════════════════════
%%
%% If you prefer brevity, delete everything below this point
%% and write a single honest paragraph. Example:
%%
%% "I am a UX researcher at [company]. I use ChatGPT to help
%%  draft interview guides and synthesize affinity diagrams.
%%  It speeds up the routine parts, but I've noticed I rely
%%  on its framing more than I'd like — my analysis increasingly
%%  mirrors its summaries. I've started blocking AI on days
%%  when I need to think slowly. The part I haven't resolved:
%%  I don't tell participants their data touched an AI."
%%
%% ═════════════════════════════════════════════════════════════
%% OPTION B — STRUCTURED FRAMEWORK
%% ═════════════════════════════════════════════════════════════

\section*{About My Practice}

\prompt{Briefly: your role, the kind of HCI work you do, and your
general relationship with AI tools.}
\response

%% ─────────────────────────────────────────────────────────────
\section*{Part A\enspace—\enspace GenAI Across the HCI Work Cycle}
%% ─────────────────────────────────────────────────────────────

\prompt{For each phase, describe how you use (or avoid) GenAI.
Skip phases that aren't part of your work.}

\paragraph{\phaseicon{\faSearch} Planning \& Literature Review}
\textit{\small Finding papers, summarizing themes, brainstorming questions.}
\response

\paragraph{\phaseicon{\faPencilRuler} Prototyping \& Design}
\textit{\small Wireframing, ``vibe coding,'' generating UI copy, assets.}
\response

\paragraph{\phaseicon{\faDatabase} Data Collection \& Synthesis}
\textit{\small Interview questions, transcript coding, persona creation.}
\response

\paragraph{\phaseicon{\faChartBar} Analysis \& Interpretation}
\textit{\small Pattern recognition, quantitative analysis, system development.}
\response

\paragraph{\phaseicon{\faBullhorn} Dissemination \& Communication}
\textit{\small Manuscript drafting, slide creation, reviewer responses.}
\response

%% ─────────────────────────────────────────────────────────────
\section*{Part B\enspace—\enspace Thematic Reflections}
%% ─────────────────────────────────────────────────────────────

\prompt{Answer the questions that feel most relevant.
You do not need to answer all of them.}

\paragraph{\phaseicon{\faUserShield} Agency and Originality}
How do you maintain status as the primary ``arbiter'' of your work?
Have you noticed shifts in professional identity when using AI?
\response

\paragraph{\phaseicon{\faHandshake} Co-Thinking and Delegation}
Which tasks have you delegated to AI, and which have you reclaimed
for deep, unassisted work? What drove those choices?
\response

\paragraph{\phaseicon{\faHeart} Empathy and Quality Integrity}
A specific moment where AI speed felt at odds with depth or care.
How did you respond?
\response

\paragraph{\phaseicon{\faFistRaised} Resistance and Re-imagining}
Have you refused to use AI for ethical or professional reasons?
What was the situation?
\response

\paragraph{\phaseicon{\faBalanceScale} Professional Policies and Boundaries}
How do institutional norms or your own values shape your AI use?
When do you choose to \textbf{STOP}?
\response

\sectionrule

%% ─────────────────────────────────────────────────────────────
\section*{What I Want to Bring to the Workshop}
%% ─────────────────────────────────────────────────────────────

\prompt{A tension, question, or practice you'd like to discuss.
This helps us form working groups around shared themes.}
\response

%% ─────────────────────────────────────────────────────────────
\bibliographystyle{ACM-Reference-Format}
% \bibliography{references}
%% The workshop proposal is pre-populated in references.bib
%% as \cite{sandhaus2026interrogating}. You may cite it if useful.
%% ─────────────────────────────────────────────────────────────

\end{document}
